\section{Paralelismo}
\subsection{Conceitos de Paralelismo}
Paralelismo é a técnica de dividir uma computação em partes que podem ser
executadas simultaneamente. Em vez de realizar todas as instruções de forma
sequencial em uma única unidade de processamento, um programa paralelo tenta
explorar múltiplas unidades de execução para reduzir tempo de resposta, aumentar
vazão ou processar entradas maiores. Essa ideia se tornou central em computação
de alto desempenho porque o aumento de frequência dos processadores encontrou
limites práticos de consumo de energia, dissipação de calor e complexidade de
projeto. Como observado por Sutter, a evolução de desempenho passou a depender
cada vez menos de processadores sequenciais mais rápidos e cada vez mais do uso
efetivo de concorrência e paralelismo \cite{sutterfree}. Essa mudança também foi
discutida por Asanovic et al., que descrevem o paralelismo como um caminho
necessário para continuar aumentando desempenho em arquiteturas modernas
\cite{asanoviclandscape}.

Uma forma clássica de classificar arquiteturas paralelas é a taxonomia de Flynn,
que organiza sistemas conforme o número de fluxos de instruções e de dados
processados simultaneamente \cite{flynntaxonomy}. Nessa taxonomia, SISD
(\textit{Single Instruction, Single Data}) representa o modelo sequencial
tradicional, SIMD (\textit{Single Instruction, Multiple Data}) representa uma
mesma instrução aplicada sobre vários dados, MISD (\textit{Multiple Instruction,
Single Data}) é uma categoria pouco comum na prática, e MIMD (\textit{Multiple
Instruction, Multiple Data}) representa sistemas nos quais diferentes unidades
executam instruções distintas sobre dados distintos. CPUs multicore se aproximam
do modelo MIMD, enquanto GPUs modernas combinam elementos SIMD/MIMD com grande
quantidade de threads.

\subsection{Processamento Multicore}

Processadores multicore integram vários núcleos de CPU em um mesmo chip. Cada
núcleo é uma unidade de processamento relativamente complexa, capaz de executar
um fluxo de instruções com baixa latência e bom desempenho em código geral. Em
comparação com arquiteturas sequenciais antigas, um processador multicore permite
que várias threads ou processos avancem simultaneamente. Esse modelo se tornou a
estratégia dominante para continuar aumentando desempenho sem depender apenas do
aumento de frequência \cite{hennessypatterson}.

Em programação multicore, uma estratégia comum é criar threads que compartilham a
memória do processo. Bibliotecas e modelos como POSIX Threads, OpenMP e C++
Threads permitem dividir laços, tarefas ou regiões independentes do programa. O
desafio principal é garantir corretude e desempenho ao mesmo tempo. 

\subsection{Processamento Manycore}

Arquiteturas manycore aumentam de forma mais agressiva o número de unidades de
execução. Em vez de poucos núcleos complexos, elas utilizam muitos núcleos mais
simples, otimizados para alta vazão. O objetivo principal não é minimizar a
latência de uma única thread, mas executar muitas threads em paralelo para
processar grandes volumes de dados. GPUs são o exemplo mais comum de arquitetura
manycore, embora o conceito também apareça em outros aceleradores e coprocessadores.

A diferença entre multicore e manycore não é apenas quantitativa. Em uma CPU
multicore, cada núcleo tenta executar código geral de forma eficiente, mesmo
quando há desvios, dependências e acesso irregular à memória. Em uma arquitetura
manycore, espera-se que o programa exponha uma grande quantidade de paralelismo.
Quando uma thread aguarda memória, outras threads podem ser escalonadas para
ocupar a unidade de execução. Dessa forma, o sistema esconde latência por meio de
concorrência massiva, em vez de depender apenas de caches grandes e controle
sofisticado \cite{kirkgpu}.

Esse modelo favorece problemas com alto paralelismo de dados e operações
semelhantes sobre muitos elementos. Simulações numéricas, processamento de
imagens, álgebra linear, aprendizado de máquina e certas etapas de processamento
de grafos podem se beneficiar desse tipo de arquitetura. Entretanto, problemas
irregulares exigem cuidado. Em grafos, por exemplo, vértices podem ter graus
muito diferentes, levando a desbalanceamento de carga. Além disso, os acessos à
memória dependem das arestas do grafo e podem ser pouco coalescidos, reduzindo a
eficiência da memória da GPU.

\subsection{GPUs}
GPUs (\textit{Graphics Processing Units}) surgiram como processadores
especializados para renderização gráfica. Com o tempo, sua arquitetura evoluiu
para executar também computação de propósito geral, área conhecida como GPGPU
(\textit{General-Purpose computing on Graphics Processing Units}). Essa evolução
foi impulsionada pela alta largura de banda de memória e pela capacidade de
executar milhares de threads leves em paralelo \cite{owensgpu, kirkgpu}.

No modelo CUDA da NVIDIA, a GPU executa kernels lançados pela CPU. Um kernel é
executado por uma grade de blocos; cada bloco contém várias threads. Threads de
um mesmo bloco podem cooperar por meio de memória compartilhada e barreiras
locais. Internamente, as threads são agrupadas em warps, que executam instruções
em estilo SIMD, mas especificamente SIMT (\textit{Single Instruction, Multiple Threads}). 
No SIMT, cada thread possui seu próprio contexto, mas threads de um mesmo warp executam a mesma
instrução ao mesmo tempo quando seguem o mesmo fluxo de controle
\cite{cudaguide}.

Essa organização tem consequências diretas para desempenho. Quando threads de um
warp seguem caminhos diferentes por causa de desvios condicionais, ocorre
divergência, e os caminhos precisam ser serializados. Quando threads acessam
posições próximas e alinhadas da memória global, a GPU consegue combinar acessos
em transações eficientes; quando os acessos são dispersos, a largura de banda é
menos aproveitada. Portanto, algoritmos eficientes em GPU precisam expor muito
paralelismo, reduzir divergência e organizar os acessos à memória sempre que
possível.

\subsection{GPUs de computação}
Existem GPUs projetadas especificamente para cargas de computação de alto desempenho, como
simulações, inteligência artificial, análise de dados e processamento científico. Essas GPUs 
geralmente oferecem mais memória, maior
largura de banda, suporte melhor à precisão dupla, recursos de virtualização,
correção de erros de memória e interconexões de alta velocidade.

Um exemplo recente desse tipo de sistema é o NVIDIA GB200 NVL72. Ele não deve ser
entendido apenas como uma GPU isolada, mas como uma arquitetura em escala de rack.
Um único rack integra 36 CPUs Grace e 72 GPUs Blackwell. A documentação da
NVIDIA descreve o GB200 Grace Blackwell Superchip como a combinação de uma CPU
Grace com duas GPUs Blackwell conectadas por NVLink-C2C.
Dentro do rack NVL72, as GPUs são interligadas por NVLink de
quinta geração e por NVLink Switch, formando um domínio NVLink com 72 GPUs. A
ideia é que as GPUs do rack consigam se comunicar com baixa latência e alta
largura de banda, aproximando o rack do comportamento de um acelerador
massivamente paralelo único \cite{nvidiagb200nvl72}.

As larguras de banda mostram a diferença entre os níveis da hierarquia. No nível
de memória local das GPUs, o GB200 NVL72 possui 13,4 TB de memória HBM3E agregada
e 576 TB/s de largura de banda agregada de memória. No nível de comunicação entre
GPUs, o domínio NVLink do rack oferece 130 TB/s de largura de banda agregada para
comunicação GPU-GPU. Já um superchip GB200, composto por uma CPU Grace e duas
GPUs Blackwell, oferece 372 GB de HBM3E, 16 TB/s de largura de banda de memória e
3,6 TB/s de largura de banda NVLink. Assim, acessos à memória local da GPU são os
mais rápidos, comunicações entre GPUs dentro do mesmo domínio NVLink ainda são de
altíssima vazão, e comunicações para fora do rack passam a depender da rede do
cluster, como InfiniBand ou Ethernet de alto desempenho \cite{nvidiagb200nvl72}.

\section{CUDA}
\label{sec:cuda}

CUDA é uma plataforma de programação paralela da NVIDIA para execução de
programas em GPUs. O modelo de programação separa a execução entre host e
dispositivo: o host, normalmente a CPU, coordena a aplicação e lança funções que
executam no dispositivo, chamadas de kernels. Esses kernels são executados por
muitas threads em paralelo, permitindo explorar o alto grau de paralelismo
disponível nas GPUs \cite{cudaguide}.

As threads em CUDA são organizadas hierarquicamente. Um kernel é lançado como uma
grade de blocos, e cada bloco contém um conjunto de threads. Threads de um mesmo
bloco podem cooperar por meio de memória compartilhada e sincronizações locais,
enquanto blocos distintos são, em geral, escalonados de forma independente. Essa
organização permite mapear problemas paralelos para diferentes granularidades de
execução: uma thread pode tratar um elemento individual, um bloco pode tratar uma
região de dados, e a grade completa pode cobrir uma entrada grande.

O desempenho em CUDA depende fortemente da forma como as threads acessam a
memória. GPUs oferecem alta largura de banda, mas essa largura de banda é melhor
aproveitada quando threads próximas acessam posições próximas da memória. Em
aplicações com acesso regular, como operações sobre vetores ou matrizes densas,
esse padrão é mais fácil de obter. Em grafos, por outro lado, a vizinhança de
cada vértice pode estar distribuída de forma irregular, e a quantidade de
trabalho por vértice pode variar muito. Por isso, algoritmos como BFS tendem a
ser limitados por acessos irregulares à memória, divergência entre threads,
contenção em operações atômicas e sincronizações entre níveis da busca.

Neste trabalho, CUDA é usada para executar em GPU as etapas paralelas da geração
do grafo, da construção de representações internas e da travessia por BFS. A
programação em CUDA fornece o mecanismo de paralelização local dentro de cada
GPU, enquanto o NVSHMEM, apresentado na Seção~\ref{sec:nvshmem}, é usado para
permitir comunicação entre elementos de processamento em execuções distribuídas.

\section{NVSHMEM}
\label{sec:nvshmem}
NVSHMEM é uma biblioteca de comunicação para GPUs NVIDIA baseada no modelo OpenSHMEM. 
O objetivo é oferecer um espaço de memória global entre GPUs, onde cada elemento de 
processamento possui uma memória local, que pode ser endereçada globalmente através de
operações de comunicação. O principal diferencial do NVSHMEM e o motivo de ser utilizado
neste trabalho é que as operações podem ser emitidas diretamente de um kernel CUDA,
sem envolver a CPU \cite{hsunvshmem}.

As unidades de processamento são divididas em PEs (Processing Elements). A biblioteca define
regiões de memória simétricas; essas regiões são alocadas coletivamente por todas as PEs.
A simetria permite que uma PE calcule o endereço lógico global correspondente a um
endereço lógico local e a um identificador de PE.

Pelas suas características, a biblioteca NVSHMEM é relevante para algoritmos irregulares.
Em aplicações densas ou estruturadas, a comunicação pode ser prevista e agrupada
em mensagens grandes. Em grafos, por outro lado, os destinos das atualizações
são descobertos durante a execução. Ao expandir uma fronteira de BFS, uma thread
pode visitar uma aresta, identificar que o vértice vizinho pertence a outro PE e
precisar atualizar o estado desse vértice remoto. Com NVSHMEM, essa atualização
pode ser expressa diretamente no kernel, usando uma operação remota ou atômica,
sem retornar primeiro ao host para empacotar mensagens \cite{hsunvshmem}.

\section{MPI}
MPI (\textit{Message Passing Interface}) é um padrão de comunicação para
programas paralelos em memória distribuída. Em um programa MPI, a execução é
organizada em processos independentes, chamados normalmente de ranks. Cada rank
possui seu próprio espaço de endereçamento e troca dados com os demais por meio
de chamadas explícitas de comunicação. Esse modelo é adequado para clusters,
onde a memória de uma máquina ou processo não pode ser acessada diretamente por
outro processo sem uma operação de comunicação definida \cite{mpistandard}.

\section{BFS}
A Busca em Largura, ou \textit{Breadth-First Search} (BFS), é um dos algoritmos fundamentais de travessia em grafos \cite{cormenalgorithms}. O algoritmo parte de um vértice inicial e, em seguida, visita os demais vértices em camadas, primeiro os que estão a uma aresta de distância, depois os que estão a duas arestas, e assim sucessivamente. Dessa forma, em um grafo não ponderado, a BFS calcula a menor distância, em número de arestas, entre o vértice inicial e todos os vértices alcançáveis a partir dele. A Busca em Largura também gera uma árvore chamada \textit{Árvore de Busca em Largura}, na qual o pai de um vértice é aquele que o descobriu durante a busca.

A implementação clássica utiliza uma fila. O vértice inicial é inserido na fila, marcado como visitado, e sua distância é definida como zero. Em seguida, enquanto a fila não estiver vazia, repete-se a seguinte operação retira-se um vértice $v$ da fila e examinam-se todos os vértices $u$ vizinhos de $v$. Caso $u$ ainda não tenha sido visitado, sua distância é calculada, ele é marcado como visitado e adicionado à fila.

Uma forma simplificada do algoritmo é apresentada a seguir:

\begin{lstlisting}
BFS(G, origem):
	para cada vertice v em G:
        pai[v] = none

    pai[origem] = origem
    fila = {origem}

    enquanto fila nao esta vazia:
		u = proximo elemento da fila
		para cada vizinho v de u:
			se pai[v] == none:
				pai[v] = u
				adiciona v na fila
\end{lstlisting}

\section{Graph500}
\label{sec:graph500-intro}

O Graph500 é um benchmark voltado à avaliação de sistemas 
de computação de alto desempenho em aplicações centradas no processamento de grafos. 
Diferentemente de benchmarks tradicionais, como o LINPACK, que enfatizam operações 
numéricas densas e desempenho em ponto flutuante, 
o Graph500 foi proposto para representar aplicações intensivas em dados. 
Nesse tipo de aplicação, o desempenho é fortemente influenciado por acesso à memória, 
comunicação entre unidades de processamento, construção de estruturas esparsas e processamento de grafos \cite{murphygraph500, graph500spec}.

A especificação do Graph500 define um conjunto de problemas sobre um grafo sintético, 
não direcionado e ponderado. Inicialmente, o grafo é gerado como uma lista de tuplas de arestas. 
A partir dessa lista, o benchmark mede a construção de uma estrutura interna de grafo e, 
em seguida, executa kernels de análise sobre essa estrutura. 
A versão da especificação considerada neste trabalho foca, em particular, 
nos benchmarks de busca (\textit{Search}) e de caminhos mínimos (\textit{Shortest Path}) \cite{graph500spec}.

O objetivo do Graph500 não é apenas medir o tempo de uma implementação isolada 
de BFS ou SSSP. O benchmark procura exercitar o sistema em um fluxo mais próximo
 de aplicações reais de grafos: geração de uma entrada irregular, construção de uma estrutura de dados adequada, 
 execução de várias buscas a partir de diferentes vértices de origem, validação dos resultados
  e cálculo de estatísticas de desempenho. Por esse motivo, o Graph500 é especialmente relevante para este trabalho.

Este capítulo apresenta os principais elementos da especificação do Graph500 necessários para contextualizar 
a implementação estudada neste trabalho. A Seção~\ref{sec:graph500-fluxo} descreve o fluxo geral de execução do benchmark. 
A Seção~\ref{sec:graph500-geracao-grafo} apresenta o processo de geração do grafo sintético utilizado pelo Graph500, incluindo os parâmetros que determinam o tamanho da instância e as propriedades estruturais produzidas pelo modelo R-MAT. 
A Seção~\ref{sec:graph500-kernel1} discute a etapa de construção da representação interna do grafo, correspondente ao Kernel 1. 
Em seguida, a Seção~\ref{sec:graph500-kernel2} apresenta a execução da busca em largura, correspondente ao Kernel 2, 
enquanto a Seção~\ref{sec:graph500-kernel3} trata, de forma complementar, do problema de caminhos mínimos executado pelo Kernel 3. 
Por fim, a Seção~\ref{sec:graph500-validacao} descreve os critérios de validação dos resultados, 
a Seção~\ref{sec:graph500-metricas} apresenta as métricas de desempenho reportadas pelo benchmark.
% a Seção~\ref{sec:graph500-criterios} discute os critérios gerais de avaliação definidos pela especificação.


\subsection{Fluxo do Benchmark}
\label{sec:graph500-fluxo}

A execução do Graph500 é organizada em etapas bem definidas. 
Primeiro, o benchmark gera uma lista de arestas. 
Em seguida, o Kernel 1 constrói uma representação de grafo a partir dessa lista. 
Depois, são selecionados até 64 vértices de origem, denominados \textit{search keys}, 
que serão utilizados nas execuções de BFS e SSSP. 
Para cada origem, o Kernel 2 executa uma BFS e 
produz uma árvore de busca em largura representada por predecessores. 
Em outro laço, o Kernel 3 executa SSSP e produz tanto predecessores quanto distâncias. 
Por fim, os resultados são validados e as métricas de desempenho são calculadas \cite{graph500spec}.

Nem todas as etapas entram no tempo reportado pelo benchmark. 
A geração da lista de arestas e a seleção das raízes não são cronometradas. 
O tempo medido inclui a construção do grafo no Kernel 1 e as execuções dos kernels de busca. 
A validação também não é incluída nas métricas de desempenho, 
mas é obrigatória para demonstrar que os resultados produzidos pela implementação são corretos.

Uma característica importante da especificação é que os Kernels 2 e 3 são executados em laços separados. 
A execução de SSSP não pode reutilizar informações computadas pela BFS, e diferentes invocações de um mesmo kernel também não podem compartilhar informações entre si. Essa regra impede que uma implementação se beneficie de conhecimento acumulado em buscas anteriores e torna a medição mais representativa da execução de buscas independentes.

\subsection{Geração do Grafo}
\label{sec:graph500-geracao-grafo}

O Graph500 utiliza um gerador sintético escalável baseado em grafos de Kronecker, semelhante ao modelo R-MAT. 
A entrada inicial do benchmark é uma lista de tuplas de arestas. 
Cada tupla contém dois vértices, \textit{StartVertex} e \textit{EndVertex}, 
e, opcionalmente, um peso. 

O principal parâmetro de entrada do benchmark é \texttt{SCALE}. O número de vértices do grafo é definido por:

\begin{equation}
N = 2^{\mathrm{SCALE}}.
\end{equation}

O número de arestas é determinado pelo parâmetro \texttt{edgefactor}, cujo valor padrão na especificação é 16:

\begin{equation}
M = \mathrm{edgefactor} \times N.
\end{equation}

Assim, ao aumentar \texttt{SCALE} em uma unidade, o número de vértices é dobrado. Mantido o mesmo valor de \texttt{edgefactor}, o número de arestas também dobra. A Tabela~\ref{tab:graph500-classes} apresenta as classes de problema definidas na especificação.

\begin{table}[H]
\centering
\begin{tabular}{l|c|c|c}
Classe & SCALE & Edge factor & Tamanho aproximado da entrada BFS \\
\hline
Toy & 26 & 16 & 0,017 TB \\
Mini & 29 & 16 & 0,137 TB \\
Small & 32 & 16 & 1,100 TB \\
Medium & 36 & 16 & 17,592 TB \\
Large & 39 & 16 & 140,738 TB \\
Huge & 42 & 16 & 1125,900 TB \\
\end{tabular}
\caption{Classes de problema do Graph500 segundo a especificação.}
\label{tab:graph500-classes}
\end{table}

O gerador R-MAT pode ser entendido como uma divisão recursiva da matriz de adjacência. Para cada aresta, a matriz é dividida em quatro quadrantes, e o gerador escolhe em qual deles a aresta será posicionada de acordo com probabilidades predefinidas. A especificação define os parâmetros iniciadores como $A = 0{,}57$, $B = 0{,}19$, $C = 0{,}19$ e $D = 1 - (A+B+C) = 0{,}05$ \cite{graph500spec}. 
Esse processo é repetido até que todos os bits dos identificadores dos vértices da aresta sejam determinados.

Essa distribuição produz grafos com estrutura irregular, contendo poucos vértices de grau muito alto e muitos vértices de grau baixo. 
Essa característica é importante para o benchmark, pois se aproxima de propriedades observadas em grafos reais e impõe desafios relacionados a balanceamento de carga, acesso à memória e comunicação. Além disso, o gerador pode produzir laços e arestas múltiplas entre os mesmos vértices. Esses elementos fazem parte da lista de entrada fornecida ao Kernel 1, ainda que possam ser tratados posteriormente pela representação interna utilizada nos kernels seguintes.

Após a geração, a especificação exige que os identificadores dos vértices sejam permutados aleatoriamente e que a lista de arestas também seja embaralhada. Esse procedimento remove localidade artificial introduzida pelo processo de geração. Sem essa etapa, uma implementação poderia explorar padrões de localidade que não necessariamente ocorreriam em grafos reais. Em implementações paralelas, a numeração dos vértices deve permanecer globalmente consistente, e a distribuição dos dados não deve introduzir localidade indevida entre processadores \cite{graph500spec}.

\subsection{Kernel 1: Construção do Grafo}
\label{sec:graph500-kernel1}

O Kernel 1 transforma a lista de tuplas de arestas em uma representação interna do grafo, que será utilizada pelos kernels posteriores. A especificação não impõe uma estrutura de dados específica. Uma implementação pode utilizar matrizes esparsas, listas de adjacência, CSR (\textit{Compressed Sparse Row}), estruturas distribuídas ou outras representações adequadas ao hardware-alvo. 
A exigência central é que a estrutura construída permita a execução correta dos kernels seguintes e a validação dos resultados produzidos \cite{graph500spec}.

Esse kernel recebe apenas a lista de arestas e o tamanho dessa lista. Informações como o número total de vértices devem ser determinadas pela própria implementação, caso sejam necessárias. O grafo resultante deve representar uma entrada não direcionada. Portanto, uma aresta entre $u$ e $v$ deve permitir travessia em ambas as direções.
O tempo de construção do grafo é incluído nas métricas do Graph500. 
Essa escolha é relevante porque, em aplicações reais, o custo de transformar uma lista de relações em uma estrutura eficiente para processamento pode ser significativo. Em ambientes paralelos, essa etapa pode envolver particionamento, redistribuição de arestas, cálculo de graus, alocação de memória e montagem das listas de adjacência.

Depois de construída, a representação do grafo não pode ser modificada entre os kernels de forma a favorecer uma execução específica. A especificação permite reservar espaço para marcações temporárias ou travas, desde que esse espaço seja utilizado apenas durante a execução do kernel corrente. Essa regra separa a estrutura permanente do grafo dos dados temporários empregados em cada busca.

Antes da execução dos kernels de busca, são selecionados aleatoriamente até 64 vértices de origem. 

\subsection{Kernel 2: Busca em Largura}
\label{sec:graph500-kernel2}

O Kernel 2 executa uma busca em largura, ou BFS, a partir de cada raiz selecionada. 
A BFS percorre o grafo por níveis: primeiro visita a origem, depois seus vizinhos, em seguida os vizinhos ainda não visitados desses vértices, e assim sucessivamente. 
Ao final, o kernel deve produzir uma árvore de busca em largura representada por um vetor de predecessores. 
O predecessor da raiz deve ser a própria raiz, e vértices não atingíveis a partir da origem devem receber predecessor igual a $-1$ \cite{graph500spec, cormenalgorithms}.

A especificação não fixa o algoritmo interno de BFS. Uma implementação pode utilizar fila, fronteiras, \textit{bitmaps}, algoritmos direcionais, particionamento distribuído ou estratégias específicas para GPU, desde que o resultado final seja uma árvore de busca em largura válida. Essa liberdade é necessária porque diferentes arquiteturas exigem diferentes estratégias de paralelização e organização dos dados.

O padrão de acesso à memória do Kernel 2 é uma das principais motivações do Graph500. Em grafos esparsos, a BFS possui acesso dependente dos dados, baixa previsibilidade e pouca computação por aresta. Além disso, grafos gerados pelo modelo R-MAT apresentam grande variação nos graus dos vértices, o que dificulta o balanceamento de carga. 
Em sistemas paralelos, múltiplas threads ou processos podem tentar descobrir o mesmo vértice simultaneamente, gerando contenção e exigindo mecanismos de sincronização.

A especificação permite uma condição de corrida benigna na escolha do predecessor de um vértice. Quando vértices do nível $k$ descobrem vértices do nível $k+1$, não é necessário garantir que o primeiro visitante seja sempre registrado como predecessor. Se mais de um vértice do nível anterior puder descobrir corretamente o mesmo vértice, qualquer um deles pode ser aceito. O requisito é que a estrutura final satisfaça as propriedades de uma árvore de busca em largura.

\subsection{Kernel 3: Caminhos Mínimos de Fonte Única}
\label{sec:graph500-kernel3}

O Kernel 3 executa SSSP (\textit{Single Source Shortest Paths}), isto é, calcula a menor distância da raiz para todos os vértices atingíveis em um grafo com pesos não negativos. A saída inclui tanto um vetor de distâncias quanto uma árvore de predecessores. Assim como na BFS, o predecessor da origem deve ser a própria origem, e vértices não atingíveis devem ter predecessor igual a $-1$ \cite{graph500spec, cormenalgorithms}.

Esse kernel é executado separadamente da BFS e não pode reutilizar dados produzidos pelo Kernel 2. Embora muitos algoritmos de SSSP também possam ser organizados em torno de fronteiras, a presença de pesos nas arestas introduz dependências adicionais e modifica a ordem em que os vértices devem ser processados. A especificação também permite corridas benignas em SSSP, desde que as distâncias e a árvore final representem caminhos mínimos corretos.

Neste trabalho, a ênfase principal está na BFS, pois ela corresponde ao Kernel 2 do Graph500 e constitui o núcleo da implementação estudada.

\subsection{Validação}
\label{sec:graph500-validacao}

Em escalas grandes, não é prático validar o resultado por comparação com uma saída padrão fixa. O grafo é gerado de forma pseudoaleatória, e diferentes implementações paralelas podem escolher predecessores distintos, mas igualmente válidos. Por esse motivo, o Graph500 utiliza uma validação baseada em propriedades estruturais do resultado \cite{graph500spec}.

No caso da BFS, a validação verifica se o vetor de predecessores representa uma árvore de busca em largura válida. Entre as propriedades verificadas estão: a raiz deve apontar para si mesma; cada vértice visitado deve possuir uma aresta real ligando-o ao seu predecessor; vértices não visitados devem estar fora da componente atingível a partir da raiz; e os níveis induzidos pela árvore devem respeitar a propriedade da busca em largura. Em particular, para qualquer aresta do grafo, os níveis de seus extremos não podem diferir por mais de uma unidade quando ambos pertencem à árvore.

Para SSSP, a validação também verifica a consistência da árvore de predecessores, mas a propriedade de níveis é substituída por propriedades de distância. As distâncias associadas aos vértices devem ser compatíveis com os pesos das arestas, e a árvore deve representar caminhos mínimos a partir da origem. A validação é executada após cada busca, mas seu tempo não é incluído nas métricas de desempenho.


\subsection{Métricas de Desempenho}
\label{sec:graph500-metricas}

A principal métrica do Graph500 é TEPS (\textit{Traversed Edges Per Second}), ou arestas percorridas por segundo. Essa métrica foi escolhida porque BFS e SSSP em grafos esparsos realizam pouca aritmética por aresta. Portanto, medir operações de ponto flutuante por segundo não representa adequadamente o principal gargalo dessas aplicações. A métrica TEPS mede a taxa com que o sistema percorre arestas e, consequentemente, reflete melhor fatores como largura de banda efetiva, latência de memória, comunicação, sincronização e balanceamento de carga \cite{graph500spec}.

Para uma execução de um kernel, a taxa é definida por:

\begin{equation}
\mathrm{TEPS} = \frac{m}{t},
\end{equation}

em que $m$ é o número de arestas percorridas na componente visitada e $t$ é o tempo medido da execução. Para arestas que não correspondem a laços, a contagem considera o número de tuplas da componente dividido por dois, pois o grafo é não direcionado. Laços são contabilizados conforme as regras definidas na especificação.

Como cada benchmark executa várias buscas, a saída deve incluir estatísticas dos tempos, das arestas percorridas e das taxas TEPS. A especificação solicita valores como mínimo, primeiro quartil, mediana, terceiro quartil, máximo, média e desvio padrão. Para TEPS, por se tratar de uma taxa, a média reportada deve ser harmônica, e não aritmética. Essa escolha reduz distorções ao combinar taxas obtidas em buscas com tempos diferentes.

Os campos de saída incluem, entre outros, \texttt{SCALE},
\texttt{edgefactor}, \texttt{NBFS}, estatísticas de tempo da BFS,
estatísticas da quantidade de arestas percorridas e estatísticas de TEPS. O
campo de tempo de construção, normalmente reportado como
\texttt{construction\_time}, registra o tempo do Kernel 1. Quando apenas um dos
kernels é executado, os campos de TEPS do outro kernel podem ser zerados,
conforme permitido pela especificação \cite{graph500spec}.
